\documentclass[12pt,a4paper]{ctexart}

\usepackage{geometry}
\usepackage{enumitem}
\usepackage{booktabs}
\usepackage{amsmath}
\usepackage{hyperref}

\geometry{left=2.5cm,right=2.5cm,top=2.5cm,bottom=2.5cm}
\setlist[itemize]{leftmargin=2em}
\setlist[enumerate]{leftmargin=2em}

\title{计算机组成原理课程内容总结}
\author{}
\date{}

\begin{document}

\maketitle

\section{课程基本信息}

本课程为《计算机组成原理》，主要围绕计算机系统的组成、运行机制以及软硬件协同工作原理展开。课程强调从系统层面理解程序执行过程，培养学生的系统分析能力和系统设计能力。

\section{参考书目}

\begin{enumerate}
    \item 袁春风等：《计算机组成与设计（基于 RISC-V 架构）》，高等教育出版社，2019 年。
    \item 胡伟武等：《计算机体系结构基础（第二版）》，机械工业出版社，2018 年。
    \item David A. Patterson 等著，王党辉等译：《计算机组成与设计：硬件/软件接口》，机械工业出版社，2015 年。
    \item 袁春风等：《计算机系统基础》，机械工业出版社，2014 年。
    \item 白中英：《计算机组成原理试题库及实现》，科学出版社，2010 年。
    \item 白中英：《计算机组成原理》，电子工业出版社，2010 年。
    \item Andrew S. Tanenbaum 等著，刘卫东等译：《计算机组成：结构化方法》，机械工业出版社，2014 年。
    \item Yale N. Patt 等著，梁阿磊等译：《计算机系统概论》，机械工业出版社，2016 年。
\end{enumerate}

\section{本学期需要掌握的主要内容}

本课程要求学生掌握计算机系统中的基本原理、机制、模型、标准以及实现方法。

\subsection{原理与机制}

需要重点理解以下内容：

\begin{itemize}
    \item 程序的自动执行机制；
    \item CPU 的基本结构与工作过程；
    \item 虚拟存储器机制；
    \item Cache 缓存机制；
    \item 页面置换算法，如 LRU、LFU；
    \item 中断机制；
    \item DMA 机制；
    \item 程序执行过程中的异常检测与处理。
\end{itemize}

\subsection{模型与标准}

需要掌握计算机系统中常见的模型与标准，例如：

\begin{itemize}
    \item 全加器与半加器；
    \item IEEE 754 浮点数标准；
    \item PCI 总线标准；
    \item USB 接口标准。
\end{itemize}

\subsection{实现方法}

课程还涉及一定的系统设计与硬件实现方法，包括：

\begin{itemize}
    \item Vivado / ISE 工具的使用；
    \item FPGA 相关设计方法；
    \item 数字逻辑电路与计算机组成部件的实现。
\end{itemize}

\section{考核方式}

课程考核形式包括课程设计和考试。

总评成绩通常由以下部分组成：

\begin{itemize}
    \item 作业；
    \item 课堂表现；
    \item 考勤；
    \item 期末考试；
    \item 课程设计。
\end{itemize}

\section{第一讲：计算机系统组成概述}

第一讲主要讨论为什么要学习《计算机组成原理》，并通过多个程序示例说明：程序运行结果不仅取决于高级语言代码本身，还与计算机底层结构、机器指令、数据表示、存储系统、操作系统以及硬件实现密切相关。

\section{什么是计算机系统}

计算机系统可以理解为硬件与软件的统一体。

\begin{itemize}
    \item 硬件是计算机系统的躯体；
    \item 软件是计算机系统的灵魂。
\end{itemize}

只有从整体系统角度理解硬件和软件之间的关系，才能真正理解程序的执行过程。

\section{为什么学习计算机组成原理}

\subsection{程序执行并不只由高级语言决定}

程序员编写的高级语言程序最终必须被转换为机器指令，才能在计算机上执行。机器指令由二进制序列组成，是计算机硬件能够直接理解并执行的形式。

因此，程序的执行结果不仅取决于高级语言语法和语义，还取决于：

\begin{itemize}
    \item 数据在机器中的表示方式；
    \item 机器指令的含义和执行过程；
    \item CPU 内部运算电路；
    \item 存储系统结构；
    \item 操作系统机制；
    \item 编译器和语言处理系统；
    \item 指令集体系结构 ISA；
    \item 微体系结构。
\end{itemize}

\section{典型程序问题分析}

\subsection{数组求和程序中的异常}

对于如下函数：

\begin{verbatim}
sum(int a[], unsigned len)
{
    int i, sum = 0;
    for (i = 0; i <= len - 1; i++)
        sum += a[i];
    return sum;
}
\end{verbatim}

当参数 \texttt{len = 0} 时，理论上返回值应该是 0，但实际执行时可能发生访存异常。

原因在于：

\begin{itemize}
    \item \texttt{len} 是无符号整数；
    \item 当 \texttt{len - 1} 时发生无符号整数下溢；
    \item 循环条件可能被错误满足；
    \item 程序访问了非法内存地址；
    \item 异常与机器指令、虚拟地址空间和异常处理机制有关。
\end{itemize}

该问题说明，理解程序执行结果需要掌握高级语言运算规则、机器指令执行、内部运算电路以及存储系统机制。

\subsection{整数乘法溢出问题}

若 \texttt{x} 和 \texttt{y} 均为 \texttt{int} 类型，当：

\[
x = 65535
\]

执行：

\[
y = x \times x
\]

可能得到：

\[
y = -131071
\]

这是因为计算机中的整数位数有限，乘法结果超过可表示范围时会发生溢出。现实世界中平方数总是非负的，但在计算机中，由于补码表示和有限字长，平方结果可能表现为负数。

该问题说明，计算机是一个有限字长的模运算系统。

\subsection{整数取负与比较问题}

对于任意 \texttt{int} 类型变量 \texttt{x} 和 \texttt{y}，表达式：

\[
(x > y) == (-x < -y)
\]

并不总是成立。

当：

\[
x = -2147483648
\]

即最小的 32 位有符号整数时，对其取负会发生溢出，因为正数范围无法表示：

\[
2147483648
\]

因此，在计算机中一些数学上看似显然成立的性质，在有限字长和补码表示下并不一定成立。

\subsection{动态内存分配中的整数溢出}

对于如下代码：

\begin{verbatim}
int copy_array(int *array, int count) {
    int i;
    int *myarray = (int *) malloc(count * sizeof(int));
    if (myarray == NULL)
        return -1;
    for (i = 0; i < count; i++)
        myarray[i] = array[i];
    return count;
}
\end{verbatim}

当 \texttt{count} 很大时，例如：

\[
count = 2^{30} + 1
\]

则：

\[
count \times sizeof(int) = 4G + 4
\]

该结果可能发生整数溢出，导致实际申请的堆空间远小于需要的空间。之后循环复制大量数据时，会破坏堆区数据。

该问题涉及：

\begin{itemize}
    \item 乘法运算及溢出；
    \item 堆空间分配；
    \item 虚拟地址空间；
    \item 存储空间映射。
\end{itemize}

\subsection{格式化输出与参数传递问题}

对于代码：

\begin{verbatim}
#include <stdio.h>

int main()
{
    double a = 10;
    printf("a = %d\n", a);
    return 0;
}
\end{verbatim}

程序使用 \texttt{\%d} 输出 \texttt{double} 类型变量，格式说明符与实际参数类型不匹配。

在 IA-32 和 x86-64 平台上，输出结果可能不同，这是因为：

\begin{itemize}
    \item \texttt{double} 采用 IEEE 754 标准表示；
    \item 不同体系结构的过程调用约定不同；
    \item IA-32 和 x86-64 的参数传递方式不同；
    \item 浮点数寄存器和整数寄存器的使用方式不同。
\end{itemize}

该例说明，程序行为不仅取决于 C 语言代码，也取决于体系结构和调用约定。

\subsection{数组越界与栈布局问题}

对于函数：

\begin{verbatim}
double fun(int i)
{
    volatile double d[1] = {3.14};
    volatile long int a[2];
    a[i] = 1073741824;
    return d[0];
}
\end{verbatim}

当 \texttt{i} 取不同值时，可能会修改栈中的不同数据，从而导致返回值异常，甚至出现存储保护错误。

该问题涉及：

\begin{itemize}
    \item 数据在机器中的表示；
    \item 过程调用机制；
    \item 栈中局部变量的布局；
    \item 数组越界访问；
    \item 存储保护机制。
\end{itemize}

\subsection{数组访问顺序与 Cache 性能问题}

两个程序功能相同，算法复杂度也相同：

\begin{verbatim}
void copyji(int src[2048][2048], int dst[2048][2048])
{
    int i, j;
    for (j = 0; j < 2048; j++)
        for (i = 0; i < 2048; i++)
            dst[i][j] = src[i][j];
}
\end{verbatim}

\begin{verbatim}
void copyij(int src[2048][2048], int dst[2048][2048])
{
    int i, j;
    for (i = 0; i < 2048; i++)
        for (j = 0; j < 2048; j++)
            dst[i][j] = src[i][j];
}
\end{verbatim}

但是它们的执行时间可能差别很大。在某些处理器上，不同访问顺序可能导致约 21 倍的性能差异。

原因在于：

\begin{itemize}
    \item C 语言二维数组按行优先存储；
    \item 按行访问具有较好的空间局部性；
    \item 按列访问容易导致 Cache 不命中；
    \item Cache 机制显著影响程序性能。
\end{itemize}

该问题说明，时间复杂度相同的程序，实际运行时间不一定相同。

\section{系统思维}

学习计算机组成原理的核心目标之一，是建立系统思维。

所谓系统思维，就是从计算机系统整体角度出发分析和解决问题。

\subsection{系统思维的基本认识}

\begin{itemize}
    \item 高级语言语句必须转换为机器指令才能执行；
    \item 机器指令是由 0 和 1 组成的二进制序列；
    \item 计算机系统是有限字长的模运算系统；
    \item 运算器本身并不知道操作数是有符号数还是无符号数；
    \item 在计算机世界中，数学规律可能因溢出而失效；
    \item 内存访问比寄存器访问慢得多；
    \item 磁盘访问比内存访问慢得多；
    \item 进程具有独立的逻辑控制流和地址空间；
    \item 过程调用通常使用栈保存参数、返回地址和局部变量；
    \item 递归调用会带来额外指令开销，并可能导致栈溢出。
\end{itemize}

\subsection{系统思维的重要性}

只有理解计算机系统的整体结构，才能：

\begin{itemize}
    \item 正确解释程序执行结果；
    \item 写出更安全、更高效的程序；
    \item 理解软硬件之间的关系；
    \item 分析系统性能瓶颈；
    \item 设计和改进计算机系统。
\end{itemize}

\section{计算机系统的抽象层次}

计算机系统由多个抽象层次组成，不同课程关注的层次不同，但这些层次之间相互联系。

程序执行结果不仅取决于算法和程序本身，还取决于：

\begin{itemize}
    \item 高级语言；
    \item 编译器；
    \item 操作系统；
    \item 指令集体系结构；
    \item 微体系结构；
    \item 数字逻辑电路；
    \item 物理硬件。
\end{itemize}

学习计算机组成原理，就是要将这些层次联系起来，从整体上理解计算机系统的运行过程。

\section{算术逻辑单元 ALU}

\subsection{ALU 的定义}

算术逻辑单元，即 Arithmetic and Logical Unit，简称 ALU，是中央处理器 CPU 的执行单元，也是 CPU 的核心组成部分。

ALU 的主要功能是执行二进制的算术运算和逻辑运算。

\subsection{ALU 的主要功能}

ALU 通常能够完成以下操作：

\begin{itemize}
    \item 加法运算；
    \item 减法运算；
    \item 与、或、非、异或等逻辑运算；
    \item 比较运算；
    \item 移位运算；
    \item 条件判断相关运算。
\end{itemize}

\subsection{数据表示形式}

在现代 CPU 体系结构中，整数数据通常采用二进制补码形式表示。补码表示使得加法和减法可以统一由加法器完成，从而简化硬件设计。

\subsection{ALU 的历史}

1945 年，数学家冯·诺依曼在 EDVAC 技术报告中首次提出了 ALU 的概念。

\subsection{简单 1 位 ALU}

简单 1 位 ALU 可以完成基本的 1 位算术与逻辑运算。多个 1 位 ALU 可以级联组成多位 ALU，例如 32 位 ALU 或 64 位 ALU。

简单 ALU 通常主要处理定点数，也就是整数运算。对于乘法和除法，简单 ALU 往往没有独立的乘除法电路模块，而是通过多次使用加法、减法和移位操作实现。

\section{课程学习的核心结论}

通过本讲内容可以得到以下结论：

\begin{enumerate}
    \item 程序执行结果不仅由高级语言代码决定，还受到机器级表示、指令执行、存储结构、操作系统和硬件机制的影响。
    \item 计算机系统是有限字长系统，整数溢出、浮点误差、地址越界等问题都可能影响程序行为。
    \item 理解程序运行结果必须从系统层面出发，而不能只停留在高级语言层面。
    \item Cache、虚拟存储器、过程调用、异常处理等机制都会影响程序的正确性和性能。
    \item 学习计算机组成原理有助于建立系统思维，为后续学习操作系统、编译原理、计算机体系结构等课程打下基础。
    \item 只有先理解系统，才能改革系统，并更好地应用和设计系统。
\end{enumerate}

\end{document}